\scanpage{70}
% archive sheet number: 71
\deleted{espace de Hilbert $H$,} soit $u$ un op. compact de $E$ dans $F$, $v$ un op. cpct de $F$ dans $G$, \deleted{posons $w=vu$. Soit $(\rho_i(w))$ la suite des valeurs propres de $\sqrt{w^*w}$} Alors, pour tout entier $n\geq1$, on a

\[
\rho_1(uv)\ldots\rho_n(uv)\leq(\rho_1(u)\rho_1(v))\ldots(\rho_n(u)\rho_n(v))\tag*{(24)}
\]

\deleted{En effet, \deleted{il est \unsure{aisé}} \deleted{$\rho_1(uv)\ldots\rho_n(uv)$} c'est évident pour $n=1$, car on a de façon générale $\rho_1(w)=\lVert(w^*w)^{1/2}\rVert=\lVert w\rVert$, et la formule se réduit donc à $\lVert uv\rVert\leq\lVert u\rVert\lVert v\rVert$. Dans le cas général, on a, puisque \deleted{\unsure{$\rho_1(uv)$}} \deleted{\unsure{ce qui}} \deleted{Et} $\rho_i(w)^2=\rho_i(w^*w)$ pour \deleted{tout} op. compact $w$ :}

\[
\begin{aligned}
\deleted{(\rho_1(uv)\ldots\rho_n(uv))^2}&\deleted{=\lambda_1((uv)(vu)^*)\ldots\lambda_n((uv)(uv)^*)=}\\
&\deleted{\lambda_1(\overset{n}{\Lambda}(uv)(uv)^*)\leq=\lambda_1}
\end{aligned}
\]

En effet, en vertu de \deleted{\unsure{form.}} (23), le premier mb est égal à $\lVert\overset{n}{\Lambda}(uv)\rVert$, i.e. $\lVert(\overset{n}{\Lambda}u)(\overset{n}{\Lambda}v)\rVert$, \deleted{(prop 1)} donc majoré par $\lVert\overset{n}{\Lambda}u\rVert\,\lVert\overset{n}{\Lambda}v\rVert$, qui n'est autre que le deuxième mb de (24) en vertu de (23).

L'idée de la démonstration de (13) et (14) sont dues à H. Weyl, ainsi que leur utilisation pour les inégalités de convexité des \unsure{fonct.|fonctions|formes}, au \deleted{\S} N\textsuperscript{o} suivant.
